\appendix

\chapter{Registered Notation}
\label{ch:registered-notation}
\label{app:notation}

This is a scoped register, not a list of global reservations.  Each row governs
only the scope named in its second column.  The same glyph may have another
standing in a disjoint scope only after that standing has been declared before
substantive use.  Within one row's scope, the registered standing and meaning
remain fixed.

\begin{longtable}{@{}p{0.14\textwidth}p{0.23\textwidth}p{0.20\textwidth}p{0.34\textwidth}@{}}
\toprule
Notation & Declared scope & Standing & Meaning \\
\midrule
\endhead
\(\Lang\) & after a language declaration & signature & Nonlogical symbols with
declared arities or sorts. \\
\(x,y,z\) & first-order syntax, semantics, and rule sheets & object-language
variables & Variables occurring in terms and formulas. \\
\(r,t\) & first-order term formation, substitution, and identity rules & term
metavariables & Members of \(\Term(\Lang)\). \\
\(\varphi,\psi,\chi\) & propositional and first-order chapters, unless a local
declaration says otherwise & formula metavariables & Members of the active
formula set. \\
\(\sigma,\tau\) & theory, model-theory, and arithmetic passages after the
sentence set is fixed & sentence metavariables & Members of
\(\Sent(\Lang)\). \\
\(\theta,\rho,\eta\) & first-order resolution chapter & simultaneous term
substitutions & Finite-support variable-to-term maps with the action and
composition declared there. \\
\(\theta_{\mathrm u},\rho_{\mathrm u},\eta_{\mathrm u}\) & automated-reasoning
chapter in this part & simultaneous term substitutions & Local visible aliases
for the preceding substitution role, used to avoid collision with a formula
metavariable in the Tseitin example. \\
\(\Gamma,\Delta,\Sigma,\Pi\) & a natural-deduction or sequent-calculus display
that declares its context representation & formula contexts & Finite sets,
multisets, or sequences exactly as declared by that active calculus. \\
\(\mathcal A_{\mathrm{in}},\mathcal A_{\mathrm{tape}}\) & Turing-machine
definitions & alphabets & Input and tape alphabets, respectively, with no
formula-context standing in that scope. \\
\(u,v,w\) & formal-language and computability passages that declare an
alphabet & strings & Finite words over the declared alphabet. \\
\(\Sigma^0_n,\Pi^0_n\) & arithmetical-hierarchy section & syntactic classes
of formulas & Classes fixed by the declared normalization and unbounded
quantifier alternation. \\
\(\boldsymbol{\Sigma}^0_n,\boldsymbol{\Pi}^0_n,
\boldsymbol{\Delta}^0_n\) & arithmetical-hierarchy section & definability
classes of relations & Classes fixed by definability in \(\mathcal N\) using
formulas from the corresponding syntactic class. \\
\(\alpha,\beta\) & cumulative-hierarchy passages & ordinal metavariables &
Ordinals in the ambient set theory. \\
\(\alpha,\beta\) & simple-type-theory passages after the local type grammar &
type metavariables & Simple types; this scope is disjoint from the preceding
ordinal scope. \\
\(T,U\) & first-order theories and metatheory after a signature is fixed &
theories & Sets of sentences in that signature. \\
\(\mathcal M,\mathcal N\) & semantic and model-theory passages & structures &
Structures for the active signature, with nonempty domains unless a local
convention says otherwise. \\
\(\abs{\mathcal M}\) & whenever \(\mathcal M\) is a structure & set & Domain
of \(\mathcal M\). \\
\(v\) & propositional and finite-matrix semantics & valuation & Map from
propositional variables to the active truth-value set; its unique
homomorphic extension \(\widehat v\) to formulas is fixed by the active
semantics. \\
\(s\) & first-order semantics & variable assignment & Map from first-order
variables to \(\abs{\mathcal M}\). \\
\(\mathcal M,s\entails\varphi\) & first-order semantics & satisfaction &
\(\varphi\) is satisfied by \(\mathcal M\) under \(s\). \\
\(\mathcal M\entails\sigma\) & sentence semantics & truth & Sentence \(\sigma\)
is true in \(\mathcal M\). \\
\(\Gamma\entails\varphi\) & after a semantic class is declared & semantic
consequence & Every interpretation and assignment in that class satisfying
\(\Gamma\) satisfies \(\varphi\). \\
\(\Gamma\derives_{\mathcal C}\varphi\) & after calculus \(\mathcal C\) is
declared & derivability & A derivation exists in \(\mathcal C\). \\
\(\varphi\equiv\psi\) & after a semantic class is declared & metalinguistic
relation & The formulas have the stated logical equivalence in that class. \\
\(\varphi\leftrightarrow\psi\) & a language containing biconditional & formula
constructor & Object-language biconditional. \\
\(a\defeq b\) & definitions in the metalanguage & metalinguistic abbreviation &
The expression on the left is defined by the expression on the right. \\
\(\varphi[t/x]\) & first-order syntax & formula operation & Capture-avoiding
substitution of \(t\) for free occurrences of \(x\) in \(\varphi\). \\
\(\FV(\varphi)\), \(\FV(\Gamma)\) & first-order syntax & set-valued operation &
Variables with a free occurrence in \(\varphi\), and their union over all
members of the formula context \(\Gamma\), respectively. \\
\(\Mod(T)\) & model theory & class & Structures satisfying every sentence of
\(T\). \\
\(\Th(\mathcal M)\) & model theory & theory & Sentences true in
\(\mathcal M\). \\
\(R^*\) & relation-closure and common-knowledge passages & relation & Reflexive
transitive closure of the locally declared relation \(R\). \\
\(\bottom,\topv\) & a formula language that declares these constants & formulas
& The bottom formula and the defined top-formula abbreviation. \\
\bottomrule
\end{longtable}

\section{Structure convention}
\label{sec:registered-notation-01}

Unless a chapter explicitly states otherwise, every first-order structure and
every sort domain is nonempty.  This convention affects quantifier laws,
Skolemization, and translations between many-sorted and one-sorted logic.

\section{Natural numbers convention}
\label{sec:registered-notation-02}

\(\N\) contains \(0\).  A finite ordinal \(n\) is
\(\set{0,\ldots,n-1}\).  The word ``positive'' therefore means greater
than \(0\), not merely a member of \(\N\).

